\chapter{Разное}

\section{Интервалы}
	\kactlimport{IntervalContainer.h}
	\kactlimport{IntervalCover.h}
	\kactlimport{ConstantIntervals.h}

\section{Тернарный поиск / золотое сечение}
	\kactlimport{GoldenSectionSearch.h}

\section{Другие алгоритмы}
	\kactlimport{FastKnapsack.h}
	\kactlimport{SimulatedAnnealing.h}

\section{Оптимизации ДП}
	\kactlimport{KnuthDP.h}
	\kactlimport{DivideAndConquerDP.h}

\section{Отладка}
\begin{itemize}[nosep,leftmargin=*]
\item \texttt{-D\_GLIBCXX\_DEBUG}: проверки итераторов и границ STL.
\item \texttt{-fsanitize=address,undefined -g}: выходы за границы,
ошибки памяти и неопределённое поведение. Использовать локально.
\item \texttt{feenableexcept(29)}: остановка на недопустимых
вещественных операциях, делении на ноль, переполнении и потере точности
при слишком малом результате (GNU/Linux).
\end{itemize}
\section{Ускорение}
\subsection{Битовые приёмы}
\begin{itemize}[nosep,leftmargin=*]
\item \texttt{x \& -x} — младший установленный бит.
\item Все подмаски m, включая нулевую:
\begin{lstlisting}
for (int s = m;; s = (s - 1) & m) {
    // обработать s
    if (s == 0)
        break;
}
\end{lstlisting}
\item Следующее число с тем же числом единиц:
\texttt{c=x\&-x; r=x+c; x=(((r\^{}x)>>2)/c)|r;}
Требуются x>0 и отсутствие переполнения r.
\item Суммы по подмаскам:
\begin{lstlisting}
for (int b = 0; b < k; ++b)
    for (int s = 0; s < (1 << k); ++s)
        if (s & (1 << b))
            d[s] += d[s ^ (1 << b)];
\end{lstlisting}
\end{itemize}
	\import{pragmas.tex}
	\kactlimport{FastMod.h}
	\kactlimport{FastInput.h}
	\kactlimport{BumpAllocator.h}
	\kactlimport{SmallPtr.h}
	\kactlimport{BumpAllocatorSTL.h}
	% \kactlimport{Unrolling.h}
	\kactlimport{SIMD.h}
